%% ============================================================================
%%  Technical Report TR-2026-03 -- The Drake Engine Adapter
%% ============================================================================

\documentclass[11pt,a4paper]{article}

\usepackage[T1]{fontenc}
\usepackage[utf8]{inputenc}
\usepackage{amsmath,amssymb}
\let\Bbbk\relax
\usepackage{newtxtext,newtxmath}
\usepackage[margin=1.0in]{geometry}
\usepackage{booktabs}
\usepackage{array}
\usepackage[expansion=false]{microtype}
\usepackage{enumitem}
\usepackage{listings}
\usepackage{xcolor}
\usepackage{colortbl}
\usepackage{titlesec}
\usepackage{fancyhdr}
\usepackage[colorlinks=true,linkcolor=linknavy,urlcolor=linknavy]{hyperref}

\definecolor{linknavy}{RGB}{20,60,110}
\definecolor{codebg}{RGB}{247,247,244}
\definecolor{coderule}{RGB}{205,205,198}
\definecolor{warnrule}{RGB}{198,150,80}

\titleformat{\section}{\normalfont\large\bfseries}{\thesection}{0.7em}{}
\titlespacing*{\section}{0pt}{1.8ex plus .4ex}{0.9ex}
\renewcommand{\arraystretch}{1.18}
\newcolumntype{L}[1]{>{\raggedright\arraybackslash}p{#1}}

\lstdefinestyle{sh}{basicstyle=\ttfamily\small,
  backgroundcolor=\color{codebg}, frame=single,
  rulecolor=\color{coderule}, framesep=6pt, breaklines=true,
  showstringspaces=false, columns=fullflexible}
\lstset{style=sh}

\newenvironment{callout}
  {\par\medskip\noindent\begin{tabular}{@{}p{0.94\textwidth}@{}}
   \arrayrulecolor{warnrule}\hline\rule{0pt}{2.6ex}}
  {\\[0.8ex]\hline\arrayrulecolor{black}\end{tabular}\par\medskip}

\newcommand{\gitsha}[1]{\texttt{#1}}

\pagestyle{fancy}
\fancyhf{}
\lhead{\footnotesize TR-2026-03}
\chead{\footnotesize The Drake Engine Adapter}
\rhead{\footnotesize \thepage}
\renewcommand{\headrulewidth}{0.4pt}

\begin{document}

\begin{center}
{\footnotesize\bfseries TECHNICAL REPORT TR-2026-03}\\[0.3em]
{\footnotesize Silent-Failure Study \quad$\cdot$\quad Week 1}\\[1.0em]
{\LARGE\bfseries The Drake Engine Adapter}\\[0.7em]
{\large Noah Parsons}\\[0.3em]
{\small 25 August 2026}
\end{center}

\vspace{0.8em}
\hrule
\vspace{1.2em}

\begin{abstract}
\noindent
This report documents the third and final engine column: Drake 1.56.0, driven
through \texttt{MultibodyPlant}. It records the installation, the modelling of
the $N$-link chain as a rigid-body mechanism, and validation against the
library-independent reference, which Drake matches to $7.2\times10^{-14}$ across
$N = 1\ldots5$. All three engines now agree. Two findings are recorded: a
coordinate-convention error of the same class as the one in TR-2026-01, caught
by the same diagnostic; and the conclusion that only one of the three
nominally portable axes transcribes directly into a rigid-body engine, which is
a scope question the study must settle before Week 2.
\end{abstract}

% ===========================================================================
\section{Installation}
% ===========================================================================

Drake was installed into a virtual environment under WSL Ubuntu 24.04.3 LTS
(Python 3.12.3, \texttt{x86\_64}):

\begin{lstlisting}
python3 -m venv ~/drake-venv
~/drake-venv/bin/pip install drake numpy scipy sympy
\end{lstlisting}

\noindent Result: \texttt{drake 1.56.0}. No \texttt{sudo} was required, since a
virtual environment sidesteps the PEP 668 restriction that Ubuntu 24.04 places
on system-wide \texttt{pip} installation.

\subsection{The execution boundary}

Drake runs on Ubuntu; the rest of the harness runs on native Windows. The
adapter is therefore written to stand alone --- it imports only \texttt{numpy},
\texttt{pydrake}, and \texttt{reference}, and emits JSON for the Windows side
to merge. This works because \texttt{reference.py} needs nothing but
\texttt{numpy} in its core path, so the same referee adjudicates Drake as
adjudicates the other two engines, on both sides of the boundary.

% ===========================================================================
\section{Modelling}
% ===========================================================================

The suite's chain is $N$ point masses on massless rods. As a Drake mechanism:

\begin{itemize}[leftmargin=1.4em,itemsep=0.25em]
  \item Body $i$'s frame origin sits at joint $i$; its point mass sits at
        $(0,0,-l)$ in that frame, at the far end of link $i$.
  \item Joint $i+1$ is placed at $(0,0,-l)$ in body $i$'s frame --- the
        position of mass $i$ --- so each link pivots about the previous mass.
  \item The revolute axis is $+y$ and gravity is $-z$, so a link hangs along
        $-z$ at zero angle, matching the suite's convention of measuring from
        the downward vertical.
\end{itemize}

\noindent Two API details are worth recording, as both cost time. Drake 1.56
has removed \texttt{SpatialInertia.PointMass}; the current equivalent is
\texttt{MakeFromCentralInertia} with a zero \texttt{RotationalInertia}. And
\texttt{RevoluteJoint} requires its two frames to be coincident, so the pivot
offset cannot be passed to the constructor --- each joint past the first needs
an explicit \texttt{FixedOffsetFrame} on the parent body.

Drake exposes the mass matrix, the bias term, and generalised gravity
separately, so the adapter solves
\begin{equation}
M(q)\,a = \tau_g(q) - C(q,v)\,v ,
\end{equation}
which carries the same content as the reference's $Ma = -(C+G)$.

% ===========================================================================
\section{A second convention trap}\label{sec:convention}
% ===========================================================================

\begin{callout}
The first validation showed Drake disagreeing with the reference by $2.5$ to
$6.4$ relative, \textbf{with the error growing in $N$}, while passing exactly
at $N=1$. This is the second occurrence of the same failure shape recorded in
TR-2026-01 \S5, and it was again a defect in the comparison rather than in the
engine.
\end{callout}

The suite's Lagrangian uses \textbf{absolute} angles measured from the downward
vertical --- its kinetic coupling $\cos(\theta_j-\theta_k)$ is only meaningful
for absolute angles. A chain of Drake revolute joints reports \textbf{relative}
angles: joint $i$ rotates its body with respect to its \emph{parent}, not with
respect to the world. The two coincide for the first link and diverge
increasingly thereafter.

The map is linear, constant, and lower triangular,
\begin{equation}
q_0 = \theta_0, \qquad q_i = \theta_i - \theta_{i-1}
\qquad\Longleftrightarrow\qquad
\theta = \operatorname{cumsum}(q),
\end{equation}
and because it is constant the same relation carries the velocities and
accelerations. Applying it on the way in and out reduced the worst mismatch
from $6.4$ to $7.2\times10^{-14}$.

\medskip
\noindent\textbf{The diagnostic generalises.} Two of the three adapters
harboured a convention mismatch invisible at the smallest system size --- the
canonical $(q,p)$ state in one case, relative joint angles in the other. In
both, the tell was identical: exact agreement at $N=1$, growing disagreement
beyond. A genuine derivation error rarely switches on at a particular system
size; a convention mismatch characteristically does, because degenerate cases
make differing conventions coincide. \emph{Any} apparent cross-engine
disagreement should be tested against this pattern before it is reported as a
finding.

% ===========================================================================
\section{Validation}
% ===========================================================================

Sixteen pseudorandom states per case, seed \texttt{20260822}, tolerance
$10^{-8}$ relative.

\begin{center}
\begin{tabular}{@{}rlr@{}}
\toprule
\textbf{$N$} & \textbf{Status} & \textbf{vs.\ reference} \\
\midrule
1 & pass & $0.00$ \\
2 & pass & $2.05\times10^{-15}$ \\
3 & pass & $1.07\times10^{-14}$ \\
4 & pass & $3.83\times10^{-14}$ \\
5 & pass & $7.21\times10^{-14}$ \\
\bottomrule
\end{tabular}
\end{center}

\noindent Drake's residual grows gently with $N$ --- roughly linearly ---
consistent with accumulation in a recursive rigid-body algorithm, against the
reference's direct closed form. At $7\times10^{-14}$ it is six orders of
magnitude inside tolerance and carries no consequence for adjudication.

% ===========================================================================
\section{Portability: only one axis transcribes}\label{sec:portability}
% ===========================================================================

Three axes were classified as portable in TR-2026-01. Building the Drake column
shows that classification was too generous.

\begin{center}
\begin{tabular}{@{}lL{0.56\textwidth}@{}}
\toprule
\textbf{Axis} & \textbf{Status in a rigid-body engine} \\
\midrule
\texttt{dof} & \textbf{Directly portable.} An $N$-link chain is a mechanism, and all three engines express it natively. \\
\texttt{mass\_ratio} & Expressible as prismatic joints with force elements, but it is a spring network rather than a mechanism. The translation is a modelling choice rather than a transcription. \\
\texttt{near\_singular} & Requires the mass matrix $\left[\begin{smallmatrix} m & cm \\ cm & m\end{smallmatrix}\right]$. A rigid-body engine yields a diagonal mass matrix for two independent prismatic degrees of freedom; the off-diagonal term can only be induced geometrically, via two prismatic axes at angle $\arccos c$ separated by a \emph{massless} intermediate body. That is expressible but contrived --- and a massless body is itself the class of degeneracy under study, so the construction would change what the axis tests. \\
\bottomrule
\end{tabular}
\end{center}

\begin{callout}
\textbf{The consequence is that \texttt{dof} may be the study's only genuine
three-way axis} --- ten of the fifty-five cases. This must be settled before
Week 2, because it determines the width of the cross-engine claim.
\end{callout}

\noindent Three responses are available, and the choice is the study's to make.
Restrict the cross-engine claim to \texttt{dof} and report the remaining axes as
single-engine characterisation. Or re-express the two spring systems on the
pendulum chain, which is what \texttt{SCOPE.md} describes in the first place
(TR-2026-01 \S7.1) and would make them mechanisms rather than abstract
Lagrangians. Or admit the contrived Drake constructions and disclose that the
translation is not faithful. The second is the most defensible and the most
work.

% ===========================================================================
\section{Status}
% ===========================================================================

\begin{center}
\begin{tabular}{@{}lccc@{}}
\toprule
 & \textbf{MechanicsDSL} & \textbf{SymPy} & \textbf{Drake} \\
\midrule
Adapter                  & native & \checkmark & \checkmark \\
Agrees with reference    & yes & yes & yes \\
Worst residual, \texttt{dof} & $3.6\times10^{-15}$ & $6.1\times10^{-15}$ & $7.2\times10^{-14}$ \\
Known scaling wall       & Hamiltonian $N{=}3$ & none at $N\leq5$ & none at $N\leq5$ \\
\bottomrule
\end{tabular}
\end{center}

\noindent\textbf{All three columns exist, and all three agree.} The Week 1
objective --- the same system expressed three ways, agreeing on a case
checkable by hand --- is met.

\medskip
\noindent\textbf{What has not been tested.} Every result so far compares
\emph{derivations}: each engine is asked for its equations of motion and
checked against the reference. None of it integrates anything. The ten timeouts
in the frozen baseline all occur during simulation, so the half of the pipeline
where the known failures live remains unexercised across engines. That is the
next measurement, and the likeliest place for a genuine disagreement to appear.

\medskip
\noindent\textbf{Artefacts.} \texttt{stress\_suite/adapter\_drake.py},
\texttt{stress\_suite/out/drake.json}. Committed at \gitsha{47b29bb}.

\medskip
\noindent\textbf{Disclosure.} The author maintains one of the engines measured.
AI assistance was used in this work. All figures were produced by executing the
committed code.

\end{document}
